\section{The cycle obstruction and the initial reductions}
\label{sec:foundations}

We first state the precise input from the companion paper. We then reduce
Theorem~\ref{thm:main} to constructing the forbidden cycle. Throughout, a
variety is integral. Boundaries are effective $\Q$-divisors, and a pair
$(X,B)$ has $X$ normal and $K_X+B$ $\Q$-Cartier. We use the usual discrepancy
definitions of log canonical (lc), Kawamata log terminal (klt), divisorial
log terminal (dlt), and terminal singularities; see \cite{KM}. Canonical
divisors on birational models are chosen compatibly. Restrictions and
intersections of $\Q$-Cartier divisors are computed using Cartier multiples.

\subsection{The obstruction supplied by the companion}

For a nef divisor $L$ over an uncountable algebraically closed field, its
\emph{nef dimension} $n(L)$ is the dimension of the target of its nef
reduction. We use only the following characterization in dimension three:
$n(L)=3$ if and only if no curve of $L$-degree zero passes through a very
general point. Here ``very general'' means outside a countable union of
proper closed subsets. These properties and existence of nef reduction are
recalled in \cite[Definition~2.1 and Theorem~2.2]{XuNuTwo}. Numerical
dimension and nef dimension need not agree.

\begin{theorem}[Cycle obstruction from the companion]
\label{thm:obstruction}
Let $k$ be an uncountable algebraically closed field of characteristic
$p>3$. There do not exist a smooth projective threefold $V$, a very ample
Cartier divisor $H_V$, and a nonzero effective Cartier divisor
$D=\sum_\alpha m_\alpha D_\alpha$ satisfying all the following conditions:
\begin{enumerate}[label=\textup{(\roman*)}]
\item The $D_\alpha$ are distinct prime divisors, $m_\alpha\in\Z_{>0}$,
the support of $D$ is
connected and has simple normal crossings, and $\gcd_\alpha m_\alpha=1$.
\item $D$ is nef, its restriction to every $D_\alpha$ is numerically
trivial, and
\[
DH_V^2>0,\qquad D^2H_V=0,\qquad n(D)=3.
\]
\item The line bundle $\mathcal O_D(D)$ on the possibly nonreduced
Cartier scheme $D$ is torsion.
\item For some integer $b\geq0$ and integers $b_\alpha$, there is an
isomorphism on a neighborhood of $\Supp D$,
\begin{equation}\label{eq:canonical-cycle}
\omega_V^{\otimes p^b}\simeq
\mathcal O_V\Bigl(\sum_\alpha b_\alpha D_\alpha\Bigr).
\end{equation}
\end{enumerate}
\end{theorem}

This is the obstruction proved in \cite[Sections~3--6]{Companion} for the
data listed in its Proposition~2.1. The end of Section~2 of that manuscript
explicitly retains only $V,D,H_V$, the torsion order, and the local canonical
relation. Under these hypotheses, its Sections~3--5 prove Corollary~5.4,
and its Proposition~6.8 gives the contradiction. Thus the obstruction
applies to any cycle with the displayed properties. In particular, the
construction below may start from a log canonical divisor with boundary.

The companion uses the torsion normal line and the local canonical
relation to control differential forms on the completion along $D$ with
poles allowed along $D$. Cartier descent then makes the Picard group of
that punctured formal neighborhood $p$-divisible. The numerical conditions
give a nonzero integer-valued degree homomorphism on the same group,
which contradicts that divisibility. We use this obstruction as a theorem;
the present proof constructs all its input data for a log canonical pair.

The canonical relation also implies an intersection identity that occurs
in the companion. Indeed, a rational section giving
\eqref{eq:canonical-cycle} has global divisor
$\sum b_\alpha D_\alpha+E$, with $\Supp E$ disjoint from $\Supp D$.
Since $D|_{D_\alpha}\equiv0$, for every Cartier divisor $A$ on $V$ we have
\begin{equation}\label{eq:canonical-intersection}
K_VDA=0.
\end{equation}

\subsection{Numerical observations}

The following elementary observation will convert zero intersection
numbers into restrictions that are numerically trivial.

\begin{lemma}\label{lem:surface-nef}
Let $T$ be an integral projective surface, $H$ an ample Cartier divisor,
and $P$ a nef $\Q$-Cartier divisor on $T$. If $PH=0$, then $P\equiv0$.
Moreover, a nonzero effective Cartier divisor on $T$ has positive
intersection with $H$.
\end{lemma}
\begin{proof}
For any integral curve $C\subset T$, a sufficiently large multiple of $H$
has an effective Cartier member $E$ containing $C$. The effective
one-cycle of $E$ has positive coefficient along $C$. Every component has
nonnegative $P$-degree, whereas $PE=0$, so $PC=0$. This proves the first
assertion. The second is positivity of the ample degree of each component
of the effective one-cycle of a nonzero Cartier divisor.
\end{proof}

In particular, if $N$ is a nonzero effective nef Cartier divisor on a
projective threefold and $N^2H=0$ for an ample Cartier divisor $H$, then
\begin{equation}\label{eq:restriction-zero}
N|_T\equiv0\qquad\text{for every prime component }T\text{ of }N.
\end{equation}
Indeed, writing $N=\sum a_TT$, the nonnegative numbers $NHT$ have
positive-coefficient sum $N^2H=0$, and Lemma~\ref{lem:surface-nef} applies.

\subsection{Reduction to a terminal underlying variety}

We use three results of Xu for projective lc threefolds in characteristic
greater than three: nonvanishing for a pseudoeffective log canonical
divisor, termination of its flips in the pseudoeffective case, and
abundance when the divisor is nef of nef dimension at most two
\cite[Theorems~1.4, 1.5, and~1.8]{XuNonvanishing}. All permit rational
boundaries. The relative MMP and crepant $\Q$-factorial dlt modifications
are available in this characteristic range, including characteristic
five \cite[Theorem~1.2]{HW}, \cite[Theorems~2.20 and~2.27]{XuNuTwo}.
We also use projective resolution of threefolds, preserving the regular
locus, and resolution of divisor supports, preserving their complement
on a regular variety; the required forms of
these results are recalled in \cite[Section~2.1]{Companion}, from
\cite{CPFirst,CP,CJS}.

\begin{proposition}\label{prop:reduction}
To prove Theorem~\ref{thm:main}, it suffices to rule out the following
data: an uncountable algebraically closed field $k$ of characteristic
$p>3$, a terminal $\Q$-factorial projective threefold $X$, an lc pair
$(X,B)$, and a nonzero effective Cartier divisor
\begin{equation}\label{eq:N0}
N_0\sim m(K_X+B),\qquad m\in\Z_{>0},
\end{equation}
such that $mK_X$ and $mB$ are Cartier, $N_0$ is nef,
$\nu(N_0)=1$, and $n(N_0)=3$.
\end{proposition}
\begin{proof}
First extend the given field to an uncountable algebraically closed
field. Integrality and normality persist, and a log resolution with its
discrepancy calculation proves that the base-changed pair remains lc.
The canonical divisor can be defined by a rational top form on the
smooth locus, so the log canonical class also base changes. Nefness
persists by ample approximation, and the prescribed intersection numbers
are unchanged.

Semiampleness descends under this extension. After making a generated
multiple also divisible by the original Cartier index, it is the base
change of a line bundle on the original variety. Global sections commute
with field extension, and faithful flatness detects surjectivity of its
evaluation map. We may therefore work over the enlarged field.

Write $L=K_X+B$ for the original log canonical class. If $n(L)\leq2$,
Xu's abundance theorem applies. Suppose $n(L)=3$. Take a crepant
projective $\Q$-factorial dlt modification $(X_d,B_d)$, and resolve $X_d$
projectively. Run a relative canonical MMP of this smooth resolution over
$X_d$. The MMP with scaling applies by \cite[Theorem~1.2]{HW}. Its output
cannot be of fibre type over the three-dimensional birational base, and
it gives a projective birational morphism
\[
g:X_t\longrightarrow X_d
\]
with $X_t$ terminal and $\Q$-factorial and $K_{X_t}$ nef over $X_d$.
Here terminality is preserved by a canonical MMP: discrepancies do not
decrease, and a divisor contracted in a negative divisorial step has
strictly positive ordinary discrepancy over the output. Relative pseudoeffectivity
causes no restriction for this birational morphism; adding a sufficiently
large pullback of an ample divisor makes the canonical class big.

The exceptional divisor $F=K_{X_t}-g^*K_{X_d}$ is $g$-nef. The negativity
lemma gives $F\leq0$. Since $X_d$ is $\Q$-factorial, the divisor
\[
B_t=g^*B_d-F
\]
is effective and satisfies
$K_{X_t}+B_t=g^*(K_{X_d}+B_d)$. Crepancy makes $(X_t,B_t)$ lc.

The new log canonical class is nef of numerical dimension one. To check
the square after pullback, observe first that $L^2T=0$ for every prime
surface $T$ on the original variety: place $T$ in an effective member of
a multiple of the ample divisor in the hypothesis, and use nonnegativity
of $L^2$ on all its components. Projection formula proves the vanishing
upstairs. Numerical nontriviality persists because a curve downstairs
can be dominated by a curve upstairs. Finally, the no-zero-degree-curve
property persists over the birational isomorphism locus. Thus the new
class still has nef dimension three.

Relabel this pair as $(X,B)$. By nonvanishing
\cite[Theorem~1.4]{XuNonvanishing}, a positive multiple of $K_X+B$ has an
effective representative $N_0$. It is nonzero because $K_X+B$ is not
numerically trivial. Increasing the multiple makes $N_0$, $mK_X$, and
$mB$ Cartier, as required.
\end{proof}

From now until the proof of Theorem~\ref{thm:main}, we work with the
data of Proposition~\ref{prop:reduction}. Fix a very ample divisor $H$ on
$X$. Then $N_0H^2>0$, $N_0^2H=0$, and
\eqref{eq:restriction-zero} holds. The next step is to establish the same
numerical triviality on the boundary components.
